\chapter{Spin Geometry and the Index Theorem}\label{spin-geometry}
\chapterstatus{planned}

The Dirac operator is a square root of the Laplacian; to define it on a manifold one needs Clifford
algebras and spin structures. Its index is computed by the Atiyah--Singer index theorem, which unifies
the Gauss--Bonnet, Hirzebruch signature and Riemann--Roch theorems.

\noindent\textbf{Prerequisites.} Chapter~\ref{multilinear-algebra} (Multilinear Algebra), Chapter~\ref{lie-groups} (Lie Groups), Chapter~\ref{characteristic-classes} (Characteristic Classes), Chapter~\ref{k-theory} (Topological K-Theory), Chapter~\ref{riemannian-geometry} (Riemannian Geometry), Chapter~\ref{elliptic-operators} (Elliptic Operators on Compact Manifolds) and Chapter~\ref{hodge-theorem} (Harmonic Forms and the Hodge Theorem).

\noindent\textbf{Planned contents.}
\begin{enumerate}
\item Clifford algebras: definition, $\mathbb{Z}/2$-grading, the classification of real and complex
Clifford algebras and their $8$-periodicity; relation with Bott periodicity.
\item The groups $\mathrm{Pin}(n)$ and $\mathrm{Spin}(n)$, the double cover of $\mathrm{SO}(n)$, spinor representations;
low-dimensional isomorphisms ($\mathrm{Spin}(3) = \mathrm{SU}(2)$, $\mathrm{Spin}(4) = \mathrm{SU}(2) \times \mathrm{SU}(2)$, \ldots).
\item Spin and $\mathrm{Spin}^c$ structures; the second Stiefel--Whitney class as obstruction; examples.
\item Spinor bundles, the spin connection, the Dirac operator; the Weitzenb\"ock formula and
Lichnerowicz's theorem; the $\hat{A}$-genus.
\item Dirac-type operators: the de Rham, signature and Dolbeault operators.
\item The Atiyah--Singer index theorem (statement and the heat-kernel sketch); the special cases of
Gauss--Bonnet--Chern, Hirzebruch's signature theorem, Hirzebruch--Riemann--Roch, and Rokhlin's theorem.
\item Supergeometry for the index theorem: Clifford modules as super vector spaces, supertraces of heat operators, Getzler's rescaling (sketch); supermanifolds, the Berezinian and the
$\hat A$-genus as a Berezin integral (survey).
\item Exercises with solutions.
\end{enumerate}
